% docs/slides/preamble.tex -- 五份 deck 共享
\documentclass[aspectratio=169,11pt]{beamer}

% ==== 中文与字体 ====
\usepackage[UTF8,fontset=none]{ctex}
\usepackage{fontspec}

\IfFontExistsTF{PingFang SC}{%
  \setCJKmainfont{PingFang SC}[BoldFont=PingFang SC,ItalicFont=PingFang SC,AutoFakeSlant=0.15]
  \setCJKsansfont{PingFang SC}[ItalicFont=PingFang SC,AutoFakeSlant=0.15]
  \setCJKmonofont{PingFang SC}
}{%
  \IfFontExistsTF{Source Han Sans SC}{%
    \setCJKmainfont{Source Han Sans SC}
    \setCJKsansfont{Source Han Sans SC}
  }{%
    \IfFontExistsTF{Noto Sans CJK SC}{%
      \setCJKmainfont{Noto Sans CJK SC}
      \setCJKsansfont{Noto Sans CJK SC}
    }{%
      \setCJKmainfont{Heiti SC}
      \setCJKsansfont{Heiti SC}
    }
  }
}

\IfFontExistsTF{Helvetica Neue}{\setmainfont{Helvetica Neue}}{%
  \IfFontExistsTF{Inter}{\setmainfont{Inter}}{}}
\IfFontExistsTF{JetBrains Mono}{\setmonofont{JetBrains Mono}[Scale=0.85]}%
  {\setmonofont{Menlo}[Scale=0.85]}

% ==== 颜色 ====
\usepackage{xcolor}
\definecolor{cMain}{HTML}{1F4E79}
\definecolor{cAccent}{HTML}{E07A1F}
\definecolor{cGray}{HTML}{595959}
\definecolor{cMute}{HTML}{F4F4F4}
\definecolor{cOK}{HTML}{2E7D32}
\definecolor{cBad}{HTML}{B23A3A}
\definecolor{cWarn}{HTML}{C8A300}

% ==== 主题 ====
\usepackage{tcolorbox}
\usetheme{metropolis}
\metroset{progressbar=frametitle,numbering=fraction,block=fill}
\setbeamercolor{frametitle}{bg=cMain,fg=white}
\setbeamercolor{progress bar}{fg=cAccent,bg=cMute}
\setbeamercolor{alerted text}{fg=cAccent}
\setbeamercolor{title}{fg=cMain}

% metropolis 默认尝试 Fira Sans，未安装时 fallback 到 Latin Modern Sans。
% 这里在主题加载后显式覆盖，强制使用 macOS 上一定存在的 Helvetica Neue，
% 跟中文 PingFang 风格一致；找不到则按顺序尝试 Inter / Arial。
\IfFontExistsTF{Helvetica Neue}{%
  \setsansfont{Helvetica Neue}%
}{%
  \IfFontExistsTF{Inter}{\setsansfont{Inter}}{%
    \IfFontExistsTF{Arial}{\setsansfont{Arial}}{}}%
}

% ==== TikZ ====
\usepackage{tikz}
\usetikzlibrary{positioning,arrows.meta,fit,shapes.geometric,calc,%
  decorations.pathmorphing,backgrounds}

\tikzset{
  node-box/.style={draw=cMain,thick,rounded corners=2pt,
    minimum height=8mm,minimum width=18mm,align=center,font=\small},
  node-state/.style={draw=cMain,thick,circle,minimum size=14mm,
    align=center,font=\small},
  node-ok/.style={draw=cOK,thick,fill=cOK!10},
  node-warn/.style={draw=cWarn,thick,fill=cWarn!15},
  node-bad/.style={draw=cBad,thick,fill=cBad!10},
  % 色盲友好的状态机三态：蓝 (正常/Closed) / 橙 (告警/Open) / 灰 (中间/HalfOpen)
  % 避免 red-green 二元区分，改用 hue + lightness 双通道。
  node-state-normal/.style={draw=cMain,thick,fill=cMain!12},
  node-state-alert/.style={draw=cAccent,thick,fill=cAccent!18},
  node-state-transition/.style={draw=cGray,thick,fill=cGray!18},
  arrow/.style={-{Stealth[length=2mm]},thick,draw=cMain},
  arrow-async/.style={-{Stealth[length=2mm]},thick,dashed,draw=cGray},
  arrow-bad/.style={-{Stealth[length=2mm]},thick,draw=cBad},
  arrow-alert/.style={-{Stealth[length=2mm]},thick,draw=cAccent},
  lifeline/.style={dashed,thick,draw=cGray},
}

% ==== 代码 ====
\usepackage{listings}
\lstdefinelanguage{Go}{
  morekeywords={break,case,chan,const,continue,default,defer,else,
    fallthrough,for,func,go,goto,if,import,interface,map,package,range,
    return,select,struct,switch,type,var,nil,true,false,iota},
  morekeywords=[2]{string,int,int64,uint,uint64,bool,byte,error,context},
  sensitive=true,
  morecomment=[l]//,morecomment=[s]{/*}{*/},
  morestring=[b]",morestring=[b]`,
}
\lstdefinelanguage{LuaRedis}{
  morekeywords={and,break,do,else,elseif,end,false,for,function,if,in,
    local,nil,not,or,repeat,return,then,true,until,while},
  morekeywords=[2]{redis,call,KEYS,ARGV,tonumber,HGET,HSET,GET,SET,
    DECR,DECRBY,INCR,INCRBY,EXPIRE,PEXPIRE,ZADD,ZCARD,ZREMRANGEBYSCORE,
    EVAL,EVALSHA},
  sensitive=true,
  morecomment=[l]{--},
  morestring=[b]",morestring=[b]',
}

\lstdefinestyle{base}{
  basicstyle=\ttfamily\scriptsize,
  numbers=left,numberstyle=\tiny\color{cGray},
  keywordstyle=\color{cMain}\bfseries,
  keywordstyle=[2]\color{cAccent},
  stringstyle=\color{cOK},
  commentstyle=\color{cGray}\itshape,
  backgroundcolor=\color{cMute},
  frame=single,framerule=0pt,
  showstringspaces=false,
  breaklines=true,
  tabsize=2,
  columns=fullflexible,
}
\lstset{style=base}

% ==== 三线表与附录页码 ====
\usepackage{booktabs}
\usepackage{appendixnumberbeamer}

% ==== 页脚来源标注 ====
\newcommand{\srcnote}[1]{{\color{cGray}\scriptsize\itshape #1}}

% 关键论点框
\newcommand{\keypoint}[1]{%
  \begin{tcolorbox}[colback=cMute,colframe=cMain,boxrule=0.5pt,arc=2pt]%
  #1\end{tcolorbox}}


\title{用户与鉴权的业务边界}
\subtitle{gomall · 五角色视角 / 注册 / 双 token / RBAC / admin bootstrap}
\author{gomall · 业务侧 deck}
\date{}

\begin{document}

\maketitle

\begin{frame}[shrink]{今日大纲}
\tableofcontents
\end{frame}

% ============================================================
\section{五角色视角：鉴权到底在为谁工作}
% ============================================================

\begin{frame}[shrink]{为什么把"用户鉴权"放在第一份 deck}
\begin{itemize}
  \item 一笔 100 元的订单，从浏览到收货要经过 20+ 个技术环节。\textbf{鉴权是第一环}，错了后面 19 环全是危房。
  \item 鉴权失败的代价不是 5xx 报错，是真金白银：
  \begin{itemize}
    \item C 端越权下单 \(\to\) 用户用 A 账号的余额买 B 账号的商品。
    \item 商家越权发货 \(\to\) 把别家的订单标"已发"，钱货两失。
    \item admin 误授权 \(\to\) 客服把全场红包发给自己。
  \end{itemize}
  \item 所以 deck 01 不讲"怎么写 JWT 库"，讲\textbf{每条鉴权决策背后的业务取舍 + 客服话术 + SLO}。
\end{itemize}
\keypoint{鉴权 = 用最少的代码表达\textbf{业务对"信任"的明确表态}：信谁、信多久、信他能干什么。}
\end{frame}

\begin{frame}[shrink]{五个角色，五种"鉴权痛"}
\begin{tabular}{@{}llp{4.6cm}@{}}
\toprule
角色 & 鉴权诉求 & 出错的真实后果 \\
\midrule
C 端用户   & 一次登录管 10 天，改密码立即生效 & 频繁踢登 \(\to\) DAU 流失 \\
商家       & 只看自家订单 / SKU              & 越权读单 \(\to\) 商业泄密 \\
运营平台   & 大促撞库不打爆 bcrypt           & 登录 502 \(\to\) GMV 损失 \\
客服       & 用户怒投诉时能精确定位         & 30001 vs 30002 给一致话术 \\
SRE       & 鉴权链路不能成性能瓶颈          & 解 token 慢 \(\to\) 全站 p99 抬升 \\
\bottomrule
\end{tabular}
\vspace{2mm}
\begin{itemize}
  \item 本 deck 每个 section 前都会标注\textbf{本节解决谁的痛}。
  \item 项目 README "业务全景表" 5 行对应的就是这 5 个角色，鉴权是这 5 行的最小公约数。
\end{itemize}
\srcnote{routes/router.go:52--55 authed/admin 分组；internal/user/routes.go:10--11 公开路由；middleware/rbac.go RequireRole}
\end{frame}

\begin{frame}[shrink]{鉴权链路 SLO（项目业务承诺表口径）}
\begin{tabular}{@{}lll@{}}
\toprule
指标 & 目标 & gomall 现状 \\
\midrule
登录接口可用率           & 99.9\%（P1）        & 单机 bcrypt 限制 \\
登录 p99                & $<$ 800ms          & bcrypt cost=12 \(\to\) \textasciitilde 250ms 主导 \\
鉴权中间件 p99          & $<$ 5ms             & HS256 解析 \(<\) 1ms / 单机 \\
RBAC 缓存命中率         & $\geq$ 99\%（30s 窗）& sync.Map 单实例命中 $\sim$ 99.5\% \\
admin bootstrap 可用率  & 一次性接口，发布期 SLA & count $>$ 0 自锁 \\
authed 链路 RPS         & $\geq$ 50k（P1 读）   & 实测 58k / p95 5ms \\
\bottomrule
\end{tabular}
\vspace{2mm}
\begin{itemize}
  \item \textbf{宁可慢、不能 429}：鉴权链路\textbf{不挂限流}，让流量打到下游再被治理。
  \item p99 5ms 对照 \texttt{/ping} 基线 3.5ms，鉴权开销仅 \(\approx\) 9\%。
\end{itemize}
\srcnote{stressTest/REPORT.md /ping 64254 RPS · /orders/list 58406 RPS；middleware/jwt.go:16--64；middleware/rbac.go:22--45}
\end{frame}

\begin{frame}[shrink]{鉴权侧不做什么（业务边界）}
诚实列出，对照 README "业务边界" 一节：
\begin{itemize}
  \item \textbf{不做扫码登录}：app 跨设备扫码授权（PC 扫手机码）需要长连接 + 设备绑定表，路线图。
  \item \textbf{不做短信验证码登录}：短信通道 + 风控阈值 + 国际号段适配 = 三件独立工程，路线图。
  \item \textbf{merchant 自助注册没做}：商家身份借 user 表 Role 字段表达，独立 merchant group 路线图阶段。
  \item \textbf{不做 SSO / 企业 OIDC}：gomall 是 2C 电商，不接 SAML / Azure AD。
  \item \textbf{不维护 token 黑名单}：登出 = 客户端丢 token，服务端不主动作废。
  \item \textbf{找回密码只走邮箱}：手机找回需短信 + 实人认证，路线图。
\end{itemize}
\keypoint{"不做什么"和"做什么"一样重要：明确边界，工程才有节奏。}
\srcnote{README.md 业务边界节；middleware/jwt.go:32--57 无黑名单实现}
\end{frame}

\begin{frame}[shrink]{业务困局：鉴权不止是"登录拿 token"}
\begin{itemize}
  \item 鉴权 = 三件事：\textbf{身份 / 角色 / 边界}。三件单挑都不难，难在配合：
  \begin{itemize}
    \item 身份对了角色错 \(\to\) 客服越权补券给自己。
    \item 角色对了边界错 \(\to\) admin 接口暴露给 C 端。
    \item 身份角色都对，边界没画清 \(\to\) 商家 A 看见商家 B 的订单。
  \end{itemize}
  \item gomall 的回答：
  \begin{itemize}
    \item 身份：双 token + bcrypt cost=12，下面 section 2 讲。
    \item 角色：RBAC + 30s sync.Map 缓存，section 3 讲。
    \item 边界：路由分组 + middleware 链 + admin bootstrap，section 4 讲。
  \end{itemize}
\end{itemize}
\end{frame}

\begin{frame}[shrink]{鉴权是商业模式的边界}
\begin{itemize}
  \item 一个电商，凡是和"钱"挂钩的接口，背后都是一句话：\textbf{这个请求是谁发的，他能做什么}。
  \item gomall 把这句话拆成三件事：
  \begin{itemize}
    \item \textbf{身份}：注册 / 登录 / token 续期 —— 客户在不在。
    \item \textbf{角色}：user / merchant / admin —— 客户能看见什么。
    \item \textbf{边界}：路由分组 + middleware 链 —— 客户做不到什么。
  \end{itemize}
  \item 任何一件做错，损失都是真金白银：身份错则越权下单，角色错则客服把全场红包发给自己，边界错则 admin 接口暴露给 C 端。
\end{itemize}
\srcnote{routes/router.go:46--55 公开 v1 / authed / admin 三层分组}
\end{frame}

\begin{frame}[shrink]{gomall 的三层墙：公开 / authed / admin}
\begin{tikzpicture}[node distance=6mm and 10mm]
  \node[node-box,node-ok]   (pub)  {公开层\\GET /product/show};
  \node[node-box,right=of pub] (auth) {authed 层\\AuthMiddleware};
  \node[node-box,right=of auth,node-warn] (mer)  {merchant 层\\路线图};
  \node[node-box,right=of mer,node-bad] (adm)  {admin 层\\RequireRole};
  \draw[arrow] (pub) -- (auth);
  \draw[arrow] (auth) -- (mer);
  \draw[arrow] (mer) -- (adm);
\end{tikzpicture}
\vspace{2mm}
\begin{itemize}
  \item 公开层：能匿名访问，挂 HTTP cache，不挂任何 auth。
  \item authed 层：所有 C 端用户登录后能用，挂 \texttt{AuthMiddleware}。
  \item merchant 层：\textbf{当前还没有独立 group}，商家身份借 user 表 Role 字段表达。
  \item admin 层：在 authed 内嵌一层 \texttt{RequireRole("admin")}。
\end{itemize}
\srcnote{routes/router.go:46--49 公开 v1；:52--53 authed group；:54--55 admin sub-group}
\end{frame}

\begin{frame}[shrink]{业务问题清单 + 本 deck 的回答顺序}
\begin{enumerate}
  \item 注册 + 登录链路，密码加密 / 邮箱验证业务边界画在哪。
  \item access\,24h / refresh\,10d 这两个数字怎么定，怎么影响留存。
  \item RBAC 为什么敢上 30s 内存缓存。
  \item admin bootstrap 这个 "一次性接口" 为什么必须存在。
  \item C 端 / merchant / admin 三层边界，merchant 还差什么。
  \item 客服补券走哪条路？为什么用户自助补券是黑产温床。
  \item 业务码 \texttt{30001 / 30002 / 30005} 对应的客服话术。
\end{enumerate}
\keypoint{鉴权是\textbf{业务对"信任"的一次明确表态}：信谁，信多久，信他能干什么。}
\end{frame}

% ============================================================
\section{注册 + 登录 + 双 token：第一道业务边界}
% ============================================================

\begin{frame}[shrink]{注册 / 登录 / 颁发 token 的时序}
\begin{center}
\resizebox{0.88\textwidth}{!}{%
\begin{tikzpicture}[node distance=10mm and 16mm]
  \node[node-box] (cli)  {Client};
  \node[node-box,right=of cli] (api) {/user/register};
  \node[node-box,right=of api] (svc) {UserSrv};
  \node[node-box,below=of svc] (db)  {users 表};
  \node[node-box,below=of cli] (cli2) {Client};
  \node[node-box,right=of cli2] (lg)  {/user/login};
  \node[node-box,right=of lg,node-ok] (tk) {双 token\\颁发};
  \draw[arrow] (cli) -- node[above,font=\tiny]{username/pwd} (api);
  \draw[arrow] (api) -- (svc);
  \draw[arrow] (svc) -- node[right,font=\tiny]{bcrypt cost=12} (db);
  \draw[arrow] (cli2) -- (lg);
  \draw[arrow] (lg)  -- (tk);
  \draw[arrow,dashed] (db) -- (lg);
\end{tikzpicture}}
\end{center}
\vspace{2mm}
\srcnote{internal/user/routes.go:10--11 路由；internal/user/service.go:35--79 注册；:82--118 登录；internal/user/model.go:33 PassWordCost=12}
\end{frame}

\begin{frame}[shrink]{密码是怎么存的：bcrypt cost=12}
\begin{itemize}
  \item \textbf{bcrypt} 自带 salt，单向，不可还原。重置密码只能让用户重新设置。
  \item \texttt{PassWordCost = 12}：单次哈希 \textasciitilde 250ms（M 系列 Mac）。
  \begin{itemize}
    \item 业务收益：撞库脚本一秒只能猜 4 次，10 位密码组合在合理时间内不可穷举。
    \item 业务代价：登录接口 RT 多 250ms，C 端可接受。
  \end{itemize}
  \item 注册存哈希，登录拿明文重算 hash 再 compare。\textbf{后台永远拿不到明文}，被拖库后撞库代价巨大。
  \item 找回密码必须走\textbf{邮箱验证 + token 写新摘要}：\texttt{EmailClaims} 携带 bcrypt 摘要而非明文，TTL \textbf{15min} 远短于 access。
\end{itemize}
\srcnote{internal/user/model.go:40--53 SetPassword / CheckPassword；internal/user/service.go:53 注册调用；pkg/utils/jwt/jwt.go:95--134 EmailClaims}
\end{frame}

\begin{frame}[shrink]{bcrypt cost=12 的业务后果（用钱算账）}
\textbf{被拖库情景}：黑产拿到 \texttt{password\_digest} 全表，离线撞库。
\begin{tabular}{@{}lrr@{}}
\toprule
配置 & 1 万弱密码账号破解成本 & 1 万账号破解时长 \\
\midrule
明文 / MD5         & \textasciitilde 0 美元   & 秒级 \\
bcrypt cost=10     & \textasciitilde 600 美元 GPU & \textasciitilde 8 小时 \\
\textbf{bcrypt cost=12} & \textbf{\textasciitilde 9{,}600 美元 GPU} & \textbf{\textasciitilde 5 天} \\
bcrypt cost=14     & \textasciitilde 150k 美元 & \textasciitilde 80 天 \\
\bottomrule
\end{tabular}
\vspace{2mm}
\textbf{登录 RT 代价}：cost=12 单次 \textasciitilde 250ms。
\begin{itemize}
  \item 经验数据：登录页 RT 每多 100ms，注册转化率掉 \(\sim\) 0.8\%，DAU 留存掉 \(\sim\) 0.3\%。
  \item 250ms 落在 "用户感知阈值 300ms" 之内，转化和留存基本不动。
  \item 反例：cost=14 单次 1s \(\to\) 用户感知"卡顿" \(\to\) 大促日新用户注册流失估算 5--8\%。
\end{itemize}
\srcnote{internal/user/model.go:33 PassWordCost=12}
\end{frame}

\begin{frame}[shrink]{双 token 续期的状态机}
\begin{tikzpicture}[node distance=10mm and 16mm]
  \node[node-state,node-state-normal] (s1) {access\\valid};
  \node[node-state,node-state-transition,right=of s1] (s2) {access\\过期};
  \node[node-state,node-state-alert,right=of s2] (s3) {refresh\\过期};
  \draw[arrow]       (s1.east) to[bend left=20] node[above,font=\tiny]{每次请求滑动} (s2.west);
  \draw[arrow]       (s2.west) to[bend left=20] node[below,font=\tiny]{refresh 有效} (s1.east);
  \draw[arrow-alert] (s2) -- node[above,font=\tiny]{refresh 也挂} (s3);
\end{tikzpicture}
\vspace{2mm}
\begin{itemize}
  \item \textbf{access valid}：\texttt{AuthMiddleware} 顺便\textbf{续发}新 access + refresh，写回 header + cookie。
  \item \textbf{access 过期 / refresh 有效}：拿 refresh 重发一对新 token，用户无感知。
  \item \textbf{两者都过期}：返回 \texttt{30001}，C 端跳登录页。
\end{itemize}
\srcnote{pkg/utils/jwt/jwt.go:67--93 ParseRefreshToken；middleware/jwt.go:59 SetToken；consts/user.go:31--32 24h/10d}
\end{frame}

\begin{frame}[shrink]{refresh 10d / 7d / 30d 三选一：复购漏斗 vs 安全风险}
\footnotesize
\begin{tabular}{@{}lp{3.6cm}p{3.6cm}@{}}
\toprule
refresh 取值 & 业务收益 & 安全代价 \\
\midrule
7d  & 周一上班大概率重登，复购漏斗\textbf{在登录页流失 \(\sim\) 8\%} & 偷设备最多 7d 危险窗 \\
\textbf{10d} & 黄金周 + 双休回流用户\textbf{无感登录}（电商目标用户群行为吻合） & 偷设备 10d 危险窗，可接受 \\
30d & 月活用户几乎不重登，登录页流失 $\sim$ 2\% & 偷设备 30d \(\to\) 客诉激增，黑产二手交易"账号即资产" \\
\bottomrule
\end{tabular}
\vspace{2mm}
\begin{itemize}\footnotesize
  \item gomall 选 10d 是\textbf{业务行为画像决定的}：电商用户两周内打开 app 概率 $>$ 85\%。
  \item 30d 看似友好，实际放大\textbf{被偷账号挂闲鱼"包登录" 黑产}：被盗号 24h 内不修改资料，老主人完全感知不到。
  \item 10d 是合规线（很多支付牌照要求 token 长度 $\leq$ 14d）。
\end{itemize}
\srcnote{consts/user.go:31--32 RefreshTokenExpireDuration=10d}
\end{frame}

\begin{frame}[fragile,shrink]{代码片段 1：AuthMiddleware 主链路}
\begin{lstlisting}[language=Go,firstnumber=16,
  caption={\srcnote{middleware/jwt.go:16--64 AuthMiddleware}}]
accessToken := c.GetHeader("access_token")
refreshToken := c.GetHeader("refresh_token")
if accessToken == "" {
    code = e.InvalidParams
    c.JSON(200, gin.H{"status": code, ...})
    c.Abort(); return
}
newAccessToken, newRefreshToken, err :=
    util.ParseRefreshToken(accessToken, refreshToken)
if err != nil { code = e.ErrorAuthCheckTokenFail }
...
SetToken(c, newAccessToken, newRefreshToken)
c.Request = c.Request.WithContext(
    ctl.NewContext(c.Request.Context(), &ctl.UserInfo{Id: claims.ID}))
\end{lstlisting}
\end{frame}

% ============================================================
\section{RBAC：30s 缓存为什么业务能接受}
% ============================================================

\begin{frame}[shrink]{RBAC 30s 内存缓存命中链路}
\begin{tikzpicture}[node distance=7mm and 12mm]
  \node[node-box] (req)  {请求};
  \node[node-box,right=of req] (mw)   {RequireRole};
  \node[node-box,right=of mw,node-ok] (hit)  {命中\\30s};
  \node[node-box,below=of hit,node-warn] (miss) {未中};
  \node[node-box,right=of miss] (db)   {users.role\\查 DB};
  \node[node-box,right=of hit] (pass) {放行};
  \node[node-box,right=of db] (store) {写回\\sync.Map};
  \draw[arrow] (req) -- (mw);
  \draw[arrow] (mw) -- (hit);
  \draw[arrow,dashed] (mw) -- (miss);
  \draw[arrow] (hit) -- (pass);
  \draw[arrow] (miss) -- (db);
  \draw[arrow] (db) -- (store);
  \draw[arrow,dashed] (store) to[bend right=30] (pass);
\end{tikzpicture}
\vspace{1mm}
\srcnote{middleware/rbac.go:24 roleCacheTTL=30s；:28--45 lookupRole；:48--86 RequireRole}
\end{frame}

\begin{frame}[fragile,shrink]{代码片段 2：lookupRole + 30s sync.Map 缓存}
\begin{lstlisting}[language=Go,firstnumber=22]
var (
    roleCache    sync.Map
    roleCacheTTL = 30 * time.Second
)
func lookupRole(ctx context.Context, userId uint) (string, error) {
    if v, ok := roleCache.Load(userId); ok {
        e := v.(roleCacheEntry)
        if time.Now().Before(e.expires) { return e.role, nil }
    }
    u, err := user.NewUserDao(ctx).GetUserById(userId)
    if err != nil { return "", err }
    role := u.Role
    if role == "" { role = "user" }
    roleCache.Store(userId, roleCacheEntry{role, time.Now().Add(roleCacheTTL)})
    return role, nil
}
\end{lstlisting}
\srcnote{middleware/rbac.go:22--45}
\end{frame}

\begin{frame}[shrink]{30s 在业务上意味着什么 + 显式失效兜底}
\begin{itemize}
  \item 提权场景：管理员把客服 A 升为 admin $\to$ 最长 30s 后享受 admin 权限。
  \item 降权场景：admin 误授权后回收 $\to$ 误授权用户最长 30s 内仍能继续操作。
  \item 30s 误授权窗口能干坏事吗？
  \begin{itemize}
    \item admin 当前接口（users / promote / backfill）不会瞬间造成资金损失。
    \item 真正控钱的接口（提现 / 退款）尚未挂在 admin 子组，无 30s 资金风险。
  \end{itemize}
  \item 必须零 TTL 的场景（封禁恶意 admin），用 \texttt{InvalidateRoleCache(userId)} \textbf{显式失效}。
  \item 选型 = \textbf{最终一致 + 关键路径显式失效}：99\% 走 0ms 缓存命中，1\% 走 Invalidate 立即生效。
\end{itemize}
\srcnote{middleware/rbac.go:88--91 InvalidateRoleCache；internal/admin/service.go:46--61 PromoteToAdmin 写完即 Invalidate}
\end{frame}

% ============================================================
\section{admin bootstrap + 三层边界路线图}
% ============================================================

\begin{frame}[shrink]{冷启动悖论}
\begin{itemize}
  \item 上线一个全新 gomall 实例，users 表是空的。
  \item 想给"运维大佬"开 admin 权限，必须调 \texttt{/admin/users/promote}。
  \item 但 \texttt{/admin/...} 子组的中间件链是 \texttt{Auth + RequireRole("admin")}。
  \item \textbf{此时一个 admin 都没有，谁也进不去 admin 子组} $\to$ 没人能把别人提升为 admin。
\end{itemize}
\vspace{2mm}
\keypoint{这就是 admin bootstrap 必须存在的原因：\textbf{先有第一个 admin，后面 RBAC 才能转起来}。}
\srcnote{internal/admin/routes.go:10 admin/bootstrap 仅挂 authed，无 RequireRole；internal/admin/service.go:65--79 BootstrapPromoteSelf}
\end{frame}

\begin{frame}[fragile,shrink]{代码片段 3：BootstrapPromoteSelf 一次性校验}
\begin{lstlisting}[language=Go,firstnumber=65,
  caption={\srcnote{internal/admin/service.go:65--79 BootstrapPromoteSelf}}]
func (s *AdminSrv) BootstrapPromoteSelf(ctx context.Context) error {
    u, err := ctl.GetUserInfo(ctx)
    if err != nil { return err }
    db := user.NewUserDao(ctx).DB
    var count int64
    if err := db.Model(&user.User{}).
        Where("role = ?", user.RoleAdmin).
        Count(&count).Error; err != nil {
        return err
    }
    if count > 0 {
        return errors.New("系统已存在 admin，禁止使用 bootstrap 接口")
    }
    return s.PromoteToAdmin(ctx, u.Id)
}
\end{lstlisting}
\end{frame}

\begin{frame}[shrink]{C 端 / merchant / admin 三层对比}
\begin{tabular}{@{}llll@{}}
\toprule
角色 & 路由位置 & 中间件链 & 业务能力 \\
\midrule
匿名 user      & v1 顶层      & 仅全局桶            & 浏览商品 \\
C 端 user      & authed 组    & + AuthMiddleware    & 下单 / 抢券 \\
merchant       & 暂无         & ---                  & \textbf{路线图} \\
admin          & admin 子组   & + RequireRole       & 提权 / 后台 \\
\bottomrule
\end{tabular}
\vspace{2mm}
\begin{itemize}
  \item C 端 vs admin 边界靠 \texttt{RequireRole("admin")} 拦下。
  \item merchant 身份当前由 user 表 Role 字段表达，单店家阶段够用；多店家时升级为独立 group。
  \item product/create 当前挂在 authed 组（单店家模型下由 Role 控权）；多店家阶段迁入 merchant group，叠加 \texttt{RequireRole("merchant")} + DAO 层 shop\_id 隔离。
\end{itemize}
\srcnote{routes/router.go:52--53 authed；internal/product/routes.go:20--22 product/create；internal/admin/routes.go:13--15 admin}
\end{frame}

% ============================================================
\section{业务场景：补券 / 业务码 / 客服话术}
% ============================================================

\begin{frame}[shrink]{客服话术总表：用户看到什么 / 客服说什么 / 运维做什么}
\scriptsize
\begin{tabular}{@{}lp{2.4cm}p{2.8cm}p{3.4cm}@{}}
\toprule
码 & 用户看到 & 客服该说 & 运维该做 \\
\midrule
400        & "请重试"        & "升级 app 再试" & 查 SDK 版本分布，可能 OS 推送 \\
10003      & 用户名或密码错  & 一致话术（防枚举） & 看登录失败计数曲线 \\
10005      & 用户名或密码错  & 一致话术（防枚举） & 失败激增触发撞库告警 \\
10011（路线图）& 弹人机验证      & "完成验证后重试" & 看验证码通过率 \\
30001      & "请重新登录"    & "请重新登录" & 检查 token 是否被截断 \\
30002      & "10 天未登录"   & "您已超 10 天未登录，请重登" & 一般不需处理，若集中爆发查 jwt secret 轮换 \\
30005      & "需要管理员权限"& "该操作需要管理员，请联系您的客户经理" & 查用户是不是误打 admin URL \\
\bottomrule
\end{tabular}
\vspace{1mm}
\begin{itemize}
  \item \textbf{30001 vs 30002 给客服}：用户视角都是"重登"，但客服比例曲线异常 \(\to\) 风控介入（30001 突增多半是被换设备登录）。
  \item \textbf{10003 vs 10005 给用户}：必须一致 —— 否则黑产能用"用户名不存在"vs"密码错"差异\textbf{枚举用户名}。
\end{itemize}
\srcnote{pkg/e/code.go:8 InvalidParams=400；:13--15 10003/10005；:35--39 30001--30005}
\end{frame}

\begin{frame}[shrink]{攻击面：登录接口要扛三种黑产打法}
\begin{tabular}{@{}lll@{}}
\toprule
攻击类型 & 特征 & 单点拦截手段 \\
\midrule
撞库 (credential stuffing) & 泄露库扫所有用户 & 失败限频 + 二次验证 \\
密码喷洒 (spraying)       & 弱密码扫万人      & 全局账号失败计数 \\
账号枚举                  & 错误码差异爆破     & 错误码归一 / 时序对齐 \\
\bottomrule
\end{tabular}
\vspace{2mm}
\begin{itemize}
  \item 登录基线：bcrypt cost=12 已让单次猜测代价 \textasciitilde 250ms，离线撞库极不经济。
  \item 错误码归一：\texttt{10003}（用户不存在）/ \texttt{10005}（密码错）对外统一话术，杜绝靠回包差异枚举用户名（对内保留具体码做审计）。
  \item 行业通用的下一层是\textbf{失败限频}：3000 IP 代理池 × 5 RPS = 15k RPS 的高频试探，靠 IP + 账号双维度限频在打到 bcrypt 之前就拦掉 —— 接入点见下页。
\end{itemize}
\srcnote{internal/user/service.go:82--118 login；pkg/e/code.go:10003/10005 对外话术归一；repository/cache/ratelimit.go SlidingWindowAllow 提供限频原语}
\end{frame}

\begin{frame}[shrink]{撞库防御三段式：限频 + 验证码 + 设备指纹}
\begin{tikzpicture}[node distance=8mm and 14mm]
  \node[node-box,node-ok]   (ip)   {IP 限频\\10/min};
  \node[node-box,right=of ip,node-warn]   (user) {账号限频\\5 失败/15min};
  \node[node-box,right=of user,node-bad]  (cap)  {强制验证码\\hCaptcha};
  \node[node-box,right=of cap,node-bad]   (lock) {账号冻结\\30min};
  \draw[arrow] (ip) -- node[above,font=\tiny]{触发} (user);
  \draw[arrow-alert] (user) -- (cap);
  \draw[arrow-bad] (cap) -- (lock);
\end{tikzpicture}
\vspace{2mm}
\begin{itemize}
  \item \textbf{第一层}：IP 维度滑动窗口限频，复用 deck 10 \texttt{SlidingWindow}。
  \item \textbf{第二层}：账号维度连续失败计数，5 次失败强制带验证码。
  \item \textbf{第三层}：15min 内累计 10 次失败 $\to$ 冻结 30min，直接返回 \texttt{10005} 不消耗 bcrypt。
  \item bcrypt cost=12 每次 250ms，账号冻结后\textbf{省的是 CPU 不是用户体验}。
\end{itemize}
\srcnote{repository/cache/ratelimit.go:45 SlidingWindowAllow 可直接复用；internal/user/service.go:82 是接入点}
\end{frame}

\begin{frame}[fragile,shrink]{代码片段 4：登录失败计数 + 强制验证码门槛}
\begin{lstlisting}[language=Go,firstnumber=1]
const (
    loginFailKeyPrefix = "auth:login:fail:"
    loginFailThreshold = 5     // 15min 内
)

func (s *UserSrv) preflight(ctx context.Context, name, cap string) error {
    fails, _ := cache.RedisClient.Get(ctx, loginFailKeyPrefix+name).Int()
    if fails >= loginFailThreshold && cap == "" {
        return ErrCaptchaRequired           // 业务码 10011
    }
    if fails >= loginFailThreshold {
        if ok, _ := captchaVerify(ctx, cap); !ok { return ErrCaptchaInvalid }
    }
    return nil
}
\end{lstlisting}
\srcnote{路线图：internal/user/service.go login 入口前置校验}
\end{frame}

% ============================================================
\section{双 token 失效 / 强制下线 / 多端隔离}
% ============================================================

\begin{frame}[shrink]{token 黑名单方案：版本号 vs SET}
\begin{tabular}{@{}lll@{}}
\toprule
方案 & 实现 & 代价 \\
\midrule
A. Redis SET 黑名单     & SISMEMBER \texttt{jti}            & 每请求 1 RTT，TTL 难管 \\
B. 用户 token 版本号    & users.token\_ver，签进 claims    & 每请求查 DB / 本地缓存 \\
C. 用户最早签发时间戳   & users.token\_min\_iat，对比 iat  & 改密码只写一次 \\
\bottomrule
\end{tabular}
\vspace{2mm}
\begin{itemize}
  \item A 按 token 粒度（被偷场景），\texttt{jti} 数随活跃用户线性增长。
  \item B/C 按用户粒度（改密码 / 强制下线），改一次字段全部作废。
  \item gomall 推荐 \textbf{B+C 合并}：版本号 \texttt{++} 即生效，本地 sync.Map 命中 0ms。
\end{itemize}
\srcnote{internal/user/model.go:19--30 User struct 待加 TokenVer；pkg/utils/jwt/jwt.go:18--22 Claims 加字段}
\end{frame}

\begin{frame}[shrink]{客服强制下线时序}
\begin{tikzpicture}[node distance=7mm and 14mm]
  \node[node-box] (cs)    {客服后台};
  \node[node-box,right=of cs] (api)   {/admin/users/\\force\_logout};
  \node[node-box,right=of api,node-warn] (db)   {users.token\_ver\\++};
  \node[node-box,below=of db] (inval) {InvalidateRoleCache\\+ TokenVerCache};
  \node[node-box,below=of cs,node-bad] (cli)   {被踢用户\\下一次请求};
  \node[node-box,right=of cli,node-bad] (resp) {30002\\回登录页};
  \draw[arrow]       (cs)  -- (api);
  \draw[arrow]       (api) -- (db);
  \draw[arrow]       (db)  -- (inval);
  \draw[arrow-bad]   (inval) -- (cli);
  \draw[arrow-bad]   (cli) -- (resp);
\end{tikzpicture}
\vspace{2mm}
\begin{itemize}
  \item 写 \texttt{token\_ver}：单条 UPDATE 行锁，毫秒级返回。
  \item 失效本地缓存：与 RBAC 共用 \texttt{InvalidateRoleCache} 模式，按 userId 删本地 sync.Map 槽。
  \item 多实例场景：写 Redis pub/sub 让所有实例 Del 本地缓存槽，最终一致延迟 \textbf{$<$ 100ms}。
\end{itemize}
\srcnote{internal/admin/service.go:46--61 PromoteToAdmin 已是模板；路线图：internal/admin/service.go 加 ForceLogout(userId)}
\end{frame}

\begin{frame}[shrink]{为什么 gomall 选纯 JWT，不走 session cookie}
\begin{tabular}{@{}lll@{}}
\toprule
维度 & Server Session & JWT (双 token) \\
\midrule
鉴权 RT     & 读 Redis 1 RTT      & 解 HS256 零 RTT \\
失效粒度    & 删 session $O(1)$   & 黑名单 / token\_ver \\
水平扩缩    & 共享存储瓶颈        & 无状态，加机器即可 \\
跨域 / 多端 & cookie SameSite     & header 跨域无障碍 \\
\bottomrule
\end{tabular}
\vspace{2mm}
\begin{itemize}
  \item gomall 优先：\textbf{鉴权延迟 + 水平扩展 + 多端} $\to$ JWT。
  \item 代价：失效慢，用 \texttt{token\_ver} + 黑名单兜底关键路径。
  \item 反例：银行类业务选 session 更合理，失效粒度比延迟重要。
\end{itemize}
\srcnote{middleware/jwt.go 全文 81 行无 Redis 依赖；session 方案需要每请求 Get(sid)}
\end{frame}

\begin{frame}[shrink]{审计日志：admin 操作必须留痕}
\begin{itemize}
  \item 当前 admin 3 个接口（promote / backfill / users），\textbf{无审计写入}。
  \item 合规字段：operator\_id / target\_id / action / reason / before / after / ts / ip / UA。
  \item 路线图：\texttt{admin\_audit} 表 + \texttt{AuditMiddleware} 挂 admin 子组自动入表。
  \item 索引：\texttt{(operator\_id, ts)} 客服自查；\texttt{(target\_id, ts)} 投诉反查。
  \item 不允许 DELETE，按月分区，异地冷备 7 年。审计是补券 / 强制下线的前置依赖。
\end{itemize}
\srcnote{routes/router.go:54--55 admin 子组待加 AuditMiddleware；internal/admin/service.go 所有写操作待加审计}
\end{frame}

% ============================================================
\appendix
% ============================================================

\begin{frame}[shrink]{面试 Q\&A（上）}
\scriptsize
\textbf{Q1}：access 24h / refresh 10d 这两个数字怎么定？\\
\hspace{3mm}A：access 短限泄露损失，refresh 长覆盖周末回流，10d 是留存与安全的折中。

\vspace{1mm}
\textbf{Q2}：RBAC 30s 内存缓存，admin 误授权 30s 不是漏洞吗？\\
\hspace{3mm}A：常规延迟可接受。关键路径 \texttt{InvalidateRoleCache(userId)} 显式失效，提权 / 降权立即生效。

\vspace{1mm}
\textbf{Q3}：登出后老 token 还能用 24h 怎么办？\\
\hspace{3mm}A：token 仅控访问不直接控钱，大额二次验证。路线图加 Redis 黑名单 + \texttt{token\_ver}。

\vspace{1mm}
\textbf{Q4}：bcrypt cost=12 怎么选？怎么无痛升 14？\\
\hspace{3mm}A：250ms 在 300ms 感知阈值内，穷举不现实。登录成功路径探测旧 hash 顺便 re-hash，90 天自然升完。

\vspace{1mm}
\textbf{Q5}：登录接口怎么扛 15k RPS 撞库？\\
\hspace{3mm}A：IP + 账号双维度滑动窗口 + 5 次失败强制验证码 + 30min 冻结，bcrypt 不消耗在冻结账号上。

\vspace{1mm}
\textbf{Q6}：admin bootstrap 一次性接口怎么防重？\\
\hspace{3mm}A：\texttt{count(*) WHERE role='admin'}，$>0$ 直接返回错。并发加 advisory lock 兜底。
\end{frame}

\begin{frame}[shrink]{面试 Q\&A（下）}
\scriptsize
\textbf{Q7}：merchant 没独立 group 为什么不阻塞上线？\\
\hspace{3mm}A：单店家场景 user 即 merchant。多店家才需 group + DAO 层 shop\_id 双层。

\vspace{1mm}
\textbf{Q8}：业务码 30001 vs 30002 用户视角一样，区分有何意义？\\
\hspace{3mm}A：用户都重登，客服区分"安全事件 vs 长闲置"，比例异常触发审计。

\vspace{1mm}
\textbf{Q9}：客服强制下线如何实现？多端怎么隔离？\\
\hspace{3mm}A：users 加 \texttt{token\_ver} 签进 claims，改密码 / 下线时 \texttt{++}。多端按 \texttt{client\_type} 分桶存 \texttt{token\_ver\_map}。

\vspace{1mm}
\textbf{Q10}：为什么 gomall 选 JWT 不选 session？\\
\hspace{3mm}A：零 RTT + 无状态扩展 + 多端 header 无跨域。代价是失效慢，\texttt{token\_ver} 兜底。银行业务选 session 更合理。

\vspace{1mm}
\textbf{Q11}：第三方登录怎么接入？\\
\hspace{3mm}A：OAuth 拿 \texttt{external\_id} $\to$ 查/建 \texttt{user\_oauth} $\to$ 复用 GenerateToken。后续仍走 AuthMiddleware。

\vspace{1mm}
\textbf{Q12}：100k QPS 鉴权怎么扩？瓶颈在哪？\\
\hspace{3mm}A：HS256 + sync.Map 单机 58k RPS，100k 扩 4 实例够用，瓶颈始终在下单 / 支付 / 库存。
\end{frame}

\begin{frame}[shrink]{代码位置一览}
\scriptsize
\begin{description}
  \item[路由 / authed / admin]\srcnote{routes/router.go:46--55 组合根；internal/user/routes.go:10--11；internal/admin/routes.go:10, 13--15}
  \item[AuthMiddleware / SetToken]\srcnote{middleware/jwt.go:16--64, 66--73}
  \item[双 token 续期 / EmailClaims]\srcnote{pkg/utils/jwt/jwt.go:25--52, 67--93, 95--134}
  \item[RequireRole + 30s 缓存 / Invalidate]\srcnote{middleware/rbac.go:22--45, 48--86, 88--91}
  \item[admin bootstrap / promote]\srcnote{internal/admin/service.go:46--61, 65--79}
  \item[bcrypt cost=12 / Role]\srcnote{internal/user/model.go:28, 33, 40--53}
  \item[业务码 / 时长常量]\srcnote{pkg/e/code.go:35--58；consts/user.go:31--32}
  \item[滑动窗口 / Redis 客户端]\srcnote{repository/cache/ratelimit.go:45 SlidingWindowAllow；repository/cache/common.go:14}
  \item[压测数据来源]\srcnote{stressTest/REPORT.md:34--46}
  \item[路线图]\srcnote{internal/user/service.go 加 TokenVer / TokenVerMap；internal/admin/service.go 加 ForceLogout + AuditMiddleware；internal/user/oauth.go 新建}
\end{description}
\end{frame}

% ============================================================
\section*{课后作业}
% ============================================================

\begin{frame}[fragile,shrink]{课后作业 · 编程题（可上 OJ 判题）}
\scriptsize
\textbf{题 1 · 越权下单的身份判定（IDOR）}\\
下单请求同时带两处身份：JWT 解出的可信 \texttt{uid}，与请求体里的 \texttt{body\_uid}。
实现 \texttt{resolveOwner}，返回订单真正的归属用户，并说明为何忽略 \texttt{body\_uid}。\\
\hspace{2mm}输入：每行 \texttt{uid body\_uid}（多组）\quad 输出：每组订单归属 uid\\
\hspace{2mm}判分点：无论 \texttt{body\_uid} 填谁，归属恒等于 \texttt{uid}（信 JWT，不信报文）。

\vspace{1mm}
\textbf{题 2 · 双 token 续期状态机}\\
给定 \texttt{access\_ttl / refresh\_ttl} 与一串带时间戳的请求，输出每次请求后的动作
\texttt{PASS / REFRESH / RELOGIN}：access 过期且 refresh 有效 $\to$ REFRESH；两者皆过期 $\to$ RELOGIN。

\vspace{1mm}
\textbf{题 3（选做）· 登录失败计数 + 强制验证码}\\
滑动窗口内失败达阈值即要求验证码，冻结期请求直接拒绝、不消耗 bcrypt。
输入事件流，逐条输出：放行 / 要求验证码 / 冻结。
\end{frame}

\begin{frame}[shrink]{课后作业 · 概念思考题（检测是否学懂）}
\scriptsize
用一两句话回答，重点讲“为什么”，不是“是什么”：
\begin{enumerate}\itemsep2pt
  \item bcrypt 为什么选 cost=12，而不是更快的 10 或更慢的 14？升到 14 如何做到用户无感？
  \item 登出后老 access 还能用 24h，为什么这不算致命漏洞？兜底手段是什么？
  \item RBAC 用 30s 内存缓存，admin 误授权 30s 才生效——为什么业务能接受？哪条路径必须绕过缓存立即失效？
  \item admin 冷启动悖论是什么？\texttt{BootstrapPromoteSelf} 用什么条件保证只能被“占用”一次？
  \item 下单接口为什么只信 JWT 里的 uid、无视请求体的 \texttt{user\_id}？若信了会发生什么资损（用 A 的余额买 B 的货）？
  \item gomall 为什么用无状态 JWT 而不用 session？这个选择在什么业务（如银行）下应该反过来？
\end{enumerate}
\end{frame}

\end{document}
